% amsthm's \qedhere in each place where it works, and proofs without it.
\documentclass{article}
\usepackage{amsmath,amsthm,enumitem}
\begin{document}
\begin{proof}[Plain]
Ends in text.
\end{proof}
\begin{proof}[Display, no qedhere]
Ends in a display.
\[ a = b \]
\end{proof}
\begin{proof}[Bracket]
Ends in a display.
\[ a = b \qedhere \]
\end{proof}
\begin{proof}[Star]
Ends in a display.
\begin{equation*} a = b \qedhere \end{equation*}
\end{proof}
\begin{proof}[Align star]
Ends in an alignment.
\begin{align*} a &= b \\ c &= d \qedhere \end{align*}
\end{proof}
\begin{proof}[Gather star]
Ends in a gather.
\begin{gather*} a = b \\ c = d \qedhere \end{gather*}
\end{proof}
\begin{proof}[Numbered]
Ends in a numbered display.
\begin{equation} a = b \qedhere \end{equation}
\end{proof}
\begin{proof}[Numbered align]
Ends in a numbered alignment.
\begin{align} a &= b \\ c &= d \qedhere \end{align}
\end{proof}
\begin{proof}[Multline star]
Ends in a multline.
\begin{multline*} a = b \\ = c \qedhere \end{multline*}
\end{proof}
\begin{proof}[Text]
Ends in text with a box here. \qedhere
\end{proof}
\begin{proof}[List]
Ends in a list.
\begin{enumerate}
\item One.
\item Two. \qedhere
\end{enumerate}
\end{proof}
\end{document}
